% =====================================================================
\section{Таймеры и конечные автоматы}
% =====================================================================
\lessonintro
  {отмерять точные интервалы аппаратным таймером и строить программу как набор состояний с понятными переходами}
  {урок 1 (\code{millis()} без \code{delay()}), урок 3 (прерывания и правила для обработчика)}
  {светодиод мигает строго по таймеру; три светодиода работают как светофор с кнопкой пешехода}
  {три светодиода (красный, жёлтый, зелёный), три резистора \SI{220}{\ohm}, кнопка}

\subsection{Сначала: будильник и расписание}

До сих пор время отмеряла сама программа: в каждом проходе \code{loop()} она смотрела на часы
\code{millis()} и решала, не пора ли что-то сделать. Это работает, пока \code{loop()} крутится быстро.
Аппаратный таймер --- это будильник внутри чипа: он считает такты кварца и в нужный момент сам прерывает
программу, чем бы она ни была занята. Вторая идея урока --- конечный автомат. Это способ записать
поведение устройства как расписание: в каком состоянии оно находится и при каком условии переходит
в следующее. Так устроены светофор, стиральная машина и меню любого прибора.

\subsection{Аппаратный таймер}

Таймер --- это счётчик, который увеличивается с заданной частотой. Когда он досчитывает до установленного
значения, срабатывает прерывание, а счётчик начинает заново. У ESP32-S3 четыре таких таймера общего
назначения.

\begin{figure}[H]
\centering
\begin{tikzpicture}[x=1cm,y=1cm,font=\sffamily\scriptsize]
  \draw[->] (0,0) -- (12.6,0) node[right]{$t$};
  \draw[->] (0,0) -- (0,2.7) node[above]{счётчик};
  \draw[dashed,esp] (0,2) -- (12.2,2) node[right,text=esp]{порог 500\,000};
  \foreach \k in {0,1,2}
    \draw[very thick,accent] (\k*4,0) -- (\k*4+4,2) -- (\k*4+4,0);
  \foreach \k in {1,2,3}{
    \draw[esp,thick,-{Stealth[length=2mm]}] (\k*4,-0.15) -- (\k*4,-0.8);
    \node[below,text=esp] at (\k*4,-0.8) {\code{onTimer()}};
    \draw (\k*4,0.06) -- (\k*4,-0.06);
  }
  \node[above] at (2,0.05) {0,5 с};
  \node[above] at (6,0.05) {0,5 с};
  \node[above] at (10,0.05) {0,5 с};
\end{tikzpicture}
\caption{Таймер считает до порога, вызывает обработчик и начинает счёт заново. Интервал не зависит от
того, чем занята основная программа}
\end{figure}

В Arduino-ESP32 3.x таймер настраивается тремя вызовами: \code{timerBegin(частота)} запускает счётчик,
\code{timerAttachInterrupt()} назначает обработчик, \code{timerAlarm(таймер, порог, повтор, число)}
задаёт порог срабатывания.

\codecap{Мигание по аппаратному таймеру}
\codeinput{l21_timer}

\begin{caution}
Обработчик таймера --- такое же прерывание, как в уроке 3. В нём нельзя вызывать \code{delay()},
\code{Serial.print()}, работать с Wi-Fi и дисплеем. Изменили переменную или вывод --- и вышли. Общие
с основной программой переменные объявляйте \code{volatile}.
\end{caution}

\begin{center}\small
\begin{tabularx}{\textwidth}{@{}l X X@{}}
\toprule
Способ отсчёта & Когда выбирать & Ограничение \\
\midrule
\code{millis()} в \code{loop()} & мигание, опрос датчиков, тайм-ауты & точность зависит от скорости \code{loop()} \\
Аппаратный таймер & строгий период: измерения, генерация сигнала & обработчик должен быть очень коротким \\
Задача FreeRTOS с \code{vTaskDelay()} & долгая работа по расписанию (урок 8) & шаг времени --- 1 мс \\
\bottomrule
\end{tabularx}
\end{center}

\subsection{Конечный автомат}

Попробуйте описать светофор словами: «горит зелёный; если пешеход нажал кнопку и прошло не меньше
пяти секунд, включить жёлтый; через две секунды --- красный\ldots». В этом описании есть всё, что нужно
программе: \textbf{состояния} (что сейчас горит), \textbf{события} (нажатие, истёкшее время) и
\textbf{переходы} между состояниями.

\begin{figure}[H]
\centering
\begin{tikzpicture}[font=\sffamily\scriptsize,
  st/.style={blk,circle,minimum size=19mm,fill=blkmem}]
  \node[st,fill=okgreen!30] (g) at (0,0) {GREEN\\ зелёный};
  \node[st,fill=warn!30] (y) at (5,0) {YELLOW\\ жёлтый};
  \node[st,fill=esp!30] (r) at (10,0) {RED\\ красный};
  \node[st,fill=warn!15] (ry) at (5,-3.4) {RED\_YELLOW\\ кр.\,+\,жёлт.};
  \draw[arr] (g) -- node[above,align=center]{кнопка нажата\\ и прошло 5 с} (y);
  \draw[arr] (y) -- node[above]{прошло 2 с} (r);
  \draw[arr] (r) |- node[pos=0.75,below]{прошло 6 с} (ry);
  \draw[arr] (ry) -| node[pos=0.25,below]{прошло 1,5 с} (g);
\end{tikzpicture}
\caption{Диаграмма состояний светофора. Кружок --- состояние, стрелка --- переход, подпись --- условие перехода}
\end{figure}

В программе состояние хранит одна переменная. Тип \code{enum class} даёт состояниям имена вместо
загадочных чисел, а оператор \code{switch} для каждого состояния проверяет только его переходы.

Светодиоды подключены к \pin{4}, \pin{6} и \pin{7} через резисторы \SI{220}{\ohm}, а кнопка пешехода
стоит между \pin{5} и GND: внешний резистор не нужен, включена внутренняя подтяжка.

\begin{figure}[H]
\centering
\begin{circuitikz}[scale=0.95,transform shape]
  \node[blk,fill=blkcpu,minimum width=2.4cm,minimum height=6.2cm] (m) at (0,0) {};
  \node[font=\sffamily\small\bfseries] at (0,3.7) {ESP32-S3};
  \foreach \y/\p/\c in {2.4/4/red, 1.2/6/yellow, 0/7/green}{
    \node[font=\sffamily\scriptsize,anchor=east] at (1.2,\y) {\pin{\p}};
    \draw[thick] (1.2,\y) to[R,l=\SI{220}{\ohm}] (4.2,\y) to[led,color=\c!70!black] (6.8,\y) -- (8.0,\y);
  }
  \node[font=\sffamily\scriptsize,anchor=east] at (1.2,-1.5) {\pin{5}};
  \draw[thick,wgreen] (1.2,-1.5) -- (3.4,-1.5) to[push button,l=S1] (6.0,-1.5) -- (8.0,-1.5);
  \node[font=\sffamily\scriptsize,anchor=east] at (1.2,-2.7) {GND};
  \draw[thick,wblack] (1.2,-2.7) -- (8.0,-2.7);
  \draw[thick,wblack] (8.0,2.4) -- (8.0,-2.7);
  \node[font=\sffamily\scriptsize,anchor=west] at (8.2,2.4) {красный};
  \node[font=\sffamily\scriptsize,anchor=west] at (8.2,1.2) {жёлтый};
  \node[font=\sffamily\scriptsize,anchor=west] at (8.2,0) {зелёный};
  \fill[wblack] (8.0,1.2) circle (1.6pt) (8.0,0) circle (1.6pt) (8.0,-1.5) circle (1.6pt);
\end{circuitikz}
\caption{Подключение светофора: три светодиода через резисторы \SI{220}{\ohm} и кнопка пешехода.
Аноды светодиодов (длинные ножки) смотрят в сторону платы, катоды идут на общую землю}
\end{figure}

Соберём светофор по шагам.

\subsubsection{Шаг 1. Три состояния по кругу}

Начнём с самого простого: зелёный, жёлтый, красный сменяют друг друга по таймеру, кнопки пока нет. Состояние
хранится в переменной \code{state}, а время входа в него --- в \code{enteredAt}.

\codecap{Шаг 1. Светофор без кнопки}
\codeinput{l21_fsm_s1}

\begin{itemize}[leftmargin=*]
  \item \code{enum class State} --- свой тип с понятными именами вместо чисел 0, 1, 2;
  \item \code{enter(s)} --- единственное место, где меняется состояние: заодно запоминается время и переключаются светодиоды;
  \item \code{switch (state)} в \code{loop()} --- для каждого состояния проверяем только его собственный переход.
\end{itemize}

\subsubsection{Шаг 2. Четвёртое состояние}

На настоящем светофоре перед зелёным на секунду горят красный и жёлтый вместе. Это новое состояние
\code{RED\_YELLOW}: чтобы добавить его, достаточно дописать имя в \code{enum}, условие в \code{enter()} и
ещё одну ветку в \code{switch}.

\codecap{Шаг 2. Состояние «красный + жёлтый»}
\codeinput{l21_fsm_s2}

\begin{itemize}[leftmargin=*]
  \item в \code{enter()} красный и жёлтый теперь горят не только в «своих» состояниях, но и в \code{RED\_YELLOW};
  \item из \code{RED} автомат переходит не сразу в зелёный, а через \code{RED\_YELLOW};
  \item остальные ветки \code{switch} не изменились --- в этом сила автомата.
\end{itemize}

\subsubsection{Шаг 3. Кнопка пешехода}

Теперь светофор должен реагировать на пешехода: зелёный держится, пока кнопку не нажали. Нажатие
запоминается во флаге \code{request} и учитывается позже, когда автомат дойдёт до нужного состояния.

\codecap{Шаг 3. Светофор с кнопкой пешехода как конечный автомат}
\codeinput{l21_fsm}

\begin{itemize}[leftmargin=*]
  \item \code{digitalRead(BTN\_PIN) == LOW} --- кнопка нажата (подтяжка внутренняя, кнопка замыкает на GND);
  \item \code{GREEN} переходит в \code{YELLOW} только при \code{request \&\& t > 5000}: зелёный горит не меньше пяти секунд;
  \item когда пешеходы прошли, в состоянии \code{RED} флаг \code{request} сбрасывается.
\end{itemize}

Обратите внимание на три приёма. Состояние меняется только в функции \code{enter()}: там же
запоминается время входа и переключаются светодиоды, поэтому забыть что-то одно невозможно. Нажатие
кнопки не обрабатывается сразу, а запоминается во флаге \code{request}. И в программе нет ни одного
\code{delay()}: в любой момент она готова реагировать на новые события.

\begin{tip}
Сначала рисуйте диаграмму на бумаге, потом пишите код. Если на диаграмме трудно объяснить какую-то
стрелку, в программе в этом месте будет ошибка.
\end{tip}

\subsection{Самостоятельный опыт: ночной режим}

Добавьте светофору пятое состояние \code{NIGHT}: жёлтый мигает раз в секунду. Переход в ночной режим
и обратно --- по команде \code{night} и \code{day} из Serial Monitor (урок 2). Сначала дорисуйте
диаграмму: из каких состояний можно попасть в \code{NIGHT} и куда из него возвращаться?

\begin{trouble}
\begin{itemize}[leftmargin=*]
  \item \textbf{Плата перезагружается после запуска таймера} --- в обработчике вызвано что-то медленное
        (\code{Serial}, \code{delay()}) или обработчик срабатывает слишком часто.
  \item \textbf{Автомат «застрял» в одном состоянии} --- время входа не обновляется: состояние меняется
        в обход функции \code{enter()}.
  \item \textbf{Кнопка срабатывает не всегда} --- нажатие проверяется только в одном состоянии; запоминайте
        его во флаге в начале \code{loop()}.
\end{itemize}
\end{trouble}

\begin{summary}
\begin{itemize}[leftmargin=*]
  \item Аппаратный таймер вызывает обработчик со строгим периодом, независимо от основной программы.
  \item Конечный автомат --- это состояния, события и переходы; в коде это \code{enum class} и \code{switch}.
  \item Состояние меняем в одном месте, события запоминаем во флагах, \code{delay()} не используем.
\end{itemize}
\end{summary}

\begin{selfcheck}
\begin{enumerate}
  \item Чем отсчёт времени аппаратным таймером отличается от проверки \code{millis()} в \code{loop()}?
  \item Что такое состояние, событие и переход? Приведите пример из светофора.
  \item Почему состояние автомата удобно менять только в одной функции?
\end{enumerate}
\end{selfcheck}

\begin{tasks}
\begin{enumerate}
\item \lv{1} Запустите таймер с периодом \SI{1}{s} и мигайте светодиодом в обработчике.
\item \lv{1} Нарисуйте граф состояний светофора: кружки --- состояния, стрелки --- переходы.
\item \lv{2} Измените период таймера так, чтобы светодиод мигал 5 раз в секунду, и проверьте по счёту.
\item \lv{2} Добавьте светофору зелёный мигающий сигнал (3 вспышки) перед жёлтым как отдельное состояние.
\item \lv{3} Добавьте кнопку для пешехода: нажатие запрашивает красный свет в конце текущего цикла.
\item \lv{4} Сделайте кодовый замок: автомат ждёт три нажатия двух кнопок в правильном порядке и зажигает зелёный светодиод; при ошибке --- красный и возврат в начало.
\end{enumerate}
\end{tasks}
